\newcommand{\cestnevyhlasenie}{
	Týmto vyhlasujem, že som predloženú \typpracerodcestneprehlasenie{} prácu vypracoval samostatne s využitím dostupnej literatúry a vlastných vedomostí pod odborným vedením vedúceho práce. Pri vypracovaní práce bola využitá umelá inteligencia \textcolor{red}{Vlož názov použitého AI modelu/modelov} na \textcolor{red}{Vlož zoznam činností, na ktoré bola použitá umelá inteligencia, napríklad jazykovú a štylistickú korektúru textu a podobne}. Všetky informačné zdroje použité pri vypracovaní tejto práce boli riadne citované a uvedené v zozname použitej literatúry v súlade s platnými predpismi.
}

\newcommand{\podakovanie}{
	Moje úprimné poďakovanie patrí školiteľovi dizertačnej práce, za odborné vedenie, cenné pripomienky a nepretržitú podporu počas jej vypracovania.
}

\newcommand{\abstraktslovencina}{
	
	\textcolor{red}{Vlož abstrakt v slovenčine}
		
	Abstrakt obsahuje informáciu o cieľoch práce, jej stručnom obsahu a v závere abstraktu sa charakterizuje splnenie cieľa, výsledky a význam celej práce. Súčasťou abstraktu je 3 - 5 kľúčových slov. Abstrakt sa píše súvisle ako jeden odsek a jeho rozsah je spravidla 100 až 500 slov.
		
	Abstrakt musí byť pre čitateľa zrozumiteľný aj bez použitia dokumentu. Pre väčšinu článkov a častí monografií stačí abstrakt s rozsahom do 250 slov. Pre poznámky a krátke správy bude dostatočná dĺžka 100 slov. Pre dlhé dokumenty, ako sú správy a dizertácie, abstrakt zvyčajne obsahuje menej ako 500 slov. Vždy, keď je to možné, používajú sa slovesá v činnom rode. Trpný rod však možno použiť na indikatívne konštatovania a na informatívne konštatovania, ak treba súčasne zdôrazniť osobu, na ktorú sa činnosť zameriava. Používa sa tretia osoba.
}

\newcommand{\klucoveslovaslovencina}{
	\textcolor{red}{Vlož kľúčové slová v slovenčine}
}

\newcommand{\abstraktanglictina}{
	\textcolor{red}{Enter the abstract in English}
		
	The abstract provides information about the objectives of the thesis, its brief content, and concludes with a description of the achievement of the objectives, the results obtained, and the overall significance of the work. The abstract also includes 3 to 5 keywords. It should be written as a single continuous paragraph and is typically between 100 and 500 words in length.
		
	The abstract must be understandable to the reader without reference to the full document. For most articles and monograph chapters, an abstract of up to 250 words is sufficient. For notes and short reports, a length of approximately 100 words is usually adequate. For longer documents, such as reports and dissertations, the abstract generally contains fewer than 500 words. Whenever possible, verbs should be used in the active voice. However, the passive voice may be used for indicative statements and for informative statements when it is necessary to emphasize the person or object affected by the action. The third person should be used throughout.
}

\newcommand{\klucoveslovaanglictina}{
	\textcolor{red}{Enter keywords in English}
}